\documentclass[11pt, a4paper]{article}

% Packages
\usepackage[utf8]{inputenc}
\usepackage[T1]{fontenc}
\usepackage{lmodern}
\usepackage[margin=0.75in]{geometry}
\usepackage{xcolor}
\usepackage{titlesec}
\usepackage{enumitem}
\usepackage{hyperref}
\usepackage{tikz}
\usepackage{tcolorbox}
\usepackage{multicol}
\usepackage{tabularx}
\usepackage{array}
\usepackage{parskip}
\usepackage{fancyhdr}
\usepackage{changepage}

% === Grayscale Color Scheme ===
\definecolor{primary}{RGB}{60, 60, 60}
\definecolor{secondary}{RGB}{80, 80, 80}
\definecolor{accent}{RGB}{240, 240, 240}
\definecolor{textgray}{RGB}{70, 70, 70}
\definecolor{lightgray}{RGB}{245, 245, 245}

% Hyperlink setup
\hypersetup{
    colorlinks=true,
    linkcolor=primary,
    urlcolor=primary,
    pdftitle={Minje Park - CV},
    pdfauthor={Minje Park}
}

% Custom section formatting
\titleformat{\section}
    {\color{primary}\Large\bfseries\scshape}
    {}{0em}
    {}
    [\color{primary}\titlerule]

\titlespacing{\section}{0pt}{8pt}{7pt}

% Subsection formatting
\titleformat{\subsection}
    {\color{secondary}\large\bfseries}
    {}{0em}
    {}

% Page style
\newcommand{\cvrulewidth}{1.5pt}
\pagestyle{fancy}
\fancyhf{}
\fancyhead[R]{\small\color{secondary} Minje Park - Curriculum Vitae}
\renewcommand{\headrulewidth}{\cvrulewidth}
\renewcommand{\footrulewidth}{\cvrulewidth}
\lfoot{\color{textgray}\small\textit{Last updated: November 2026}}
\cfoot{\small - \thepage\ -}

% Custom commands
\newcommand{\cvtag}[1]{%
    \tikz[baseline=(tag.base)]\node[draw=primary, fill=accent, rounded corners=3pt, inner sep=3pt, text=primary] (tag) {\small #1};%
}

\newcommand{\experienceitem}[4]{
    \begin{tcolorbox}[
        colback=lightgray,
        colframe=primary,
        boxrule=0pt,
        leftrule=3pt,
        arc=0pt,
        left=8pt,
        right=8pt,
        top=5pt,
        bottom=5pt
    ]
    \textbf{#1}\hfill{\color{primary}\small #2}\par
    \vspace{1pt}
    {\small\itshape\color{secondary}#3}\par
    \vspace{3pt}
    #4
    \end{tcolorbox}
    \vspace{1pt}
}

\newcommand{\conferenceitem}[3]{%
    \noindent
    \parbox[t]{0.15\textwidth}{{\color{primary}\textbf{#1}}}%
    \hfill
    \parbox[t]{0.83\textwidth}{\textbf{#2}\\[2pt]\textit{#3}}%
    \vspace{8pt}\par
}

\newcommand{\awarditem}[3]{
    \begin{tcolorbox}[
        colback=white,
        colframe=primary,
        boxrule=0pt,
        leftrule=3pt,
        arc=0pt,
        left=5pt,
        right=5pt,
        top=4pt,
        bottom=4pt
    ]
    \textbf{#1} \\[2pt]
    #2 \\[2pt]
    {\small\textit{#3}}
    \end{tcolorbox}
    \vspace{3pt}
}

% Compact bullets right under a line of text (cancels the parskip gap)
\newenvironment{tightbullets}{\vspace{-\parskip}\begin{itemize}[leftmargin=1.2em, itemsep=2.5pt, parsep=0pt, topsep=2pt, partopsep=0pt]}{\end{itemize}}

% Section bodies sit slightly inside their headings
\newenvironment{secbody}{\begin{adjustwidth}{1em}{0pt}}{\end{adjustwidth}}

\begin{document}

% === Header ===
\vspace{-6pt}

{\Huge\bfseries\color{primary} MINJE PARK}

\vspace{6pt}
\noindent
{\small
\renewcommand{\arraystretch}{1.1}
\begin{tabularx}{\textwidth}{@{}p{0.13\textwidth}X@{}}
Email &
\href{mailto:pmj0324@snu.ac.kr}{pmj0324@snu.ac.kr} \\
Homepage &
\href{https://pmj0324.github.io/}{https://pmj0324.github.io/} \\
GitHub &
\href{https://github.com/pmj0324}{github.com/pmj0324} \\
Scholar &
\href{https://scholar.google.com/citations?user=oC5mE0AAAAAJ}{Google Scholar (Minje Park)} \\
\end{tabularx}}

\vspace{2pt}

% === Research Interests ===
\section{Research Interests}
Computational cosmology and machine learning, with a focus on:
\begin{tightbullets}
    \item Generative models (flow matching and diffusion) for cosmological simulations.
    \item Simulation-based inference (SBI) of cosmological parameters.
    \item Machine learning algorithms for scientific data, including new training objectives and architectures.
\end{tightbullets}

% === Education ===
\section{Education}
\begin{secbody}

\vspace{3.5pt}
\noindent
{\large\textbf{Seoul National University}} \hfill \textit{Mar 2027 -- (Admitted)}
\begin{tightbullets}
    \item Integrated M.S./Ph.D. Program in Physics, Department of Physics and Astronomy
\end{tightbullets}

\vspace{2pt}
\noindent
{\large\textbf{Sungkyunkwan University}} \hfill \textit{Feb 2021 -- Feb 2027 (Expected)}
\begin{tightbullets}
    \item B.S. in Physics, Double Major in Electronic \& Electrical Engineering
\end{tightbullets}
\end{secbody}

% === Research Experience ===
\section{Research Experience}
\begin{secbody}

\newcommand{\expinst}[1]{{\upshape\color{primary}#1}}
\newenvironment{expbullets}{\begin{itemize}[leftmargin=1.2em, itemsep=1.5pt, topsep=0pt, parsep=0pt]}{\end{itemize}}

\experienceitem{Seoul National University}{Jun 2026 -- Present}{\expinst{Research Intern} \textperiodcentered{} Advisor: Prof.\ Ji-Hoon Kim, Department of Physics and Astronomy}{
    \begin{expbullets}
        \item Developing generative AI models for cosmological simulation and super-resolution.
    \end{expbullets}
}

\experienceitem{University of Seoul}{Jun 2024 -- Feb 2026}{\expinst{Research Intern} \textperiodcentered{} Advisor: Prof.\ Ian James Watson, Department of Physics}{
    \begin{expbullets}
        \item Investigated position embedding strategies for a Transformer-based primary-particle reconstruction model in SWGO.
        \item Developed flow matching models to synthesize SWGO detector signals from simulation data.
    \end{expbullets}
}

\experienceitem{Sungkyunkwan University}{May 2024 -- Feb 2026}{\expinst{Research Intern} \textperiodcentered{} Advisor: Prof.\ Chang Dong Rho, Department of Physics}{
    \begin{expbullets}
        \item Built flow and diffusion models to generate neutrino events in the IceCube Observatory.
        \item Developed diffusion models and Mixture Density Networks to synthesize PMT signals for the SWGO tank simulation.
    \end{expbullets}
}

\experienceitem{Sungkyunkwan University \& University of Utah}{Aug 2023 -- Sep 2025}{\expinst{Research Intern} \textperiodcentered{} Advisor: Prof.\ Carsten Rott, Department of Physics}{
    \begin{expbullets}
        \item Developed CNN and Vision Transformer models for the IceCube Upgrade Camera System.
        \item Results presented as a poster and proceedings paper at ICRC 2025 (arXiv:2507.18525).
    \end{expbullets}
}
\end{secbody}

% === Publications ===
\section{Publications}
\begin{secbody}

\newcommand{\publication}[3]{
    \noindent
    \textbf{#1} \\
    \textit{#2} \\
    {\small #3}
    \vspace{6pt}
}

\publication{Performance Study of the IceCube Upgrade Camera System}
{C. Rott, \textbf{\textit{M. Park}}, M. Jansson, G. Iverson, S. Choi (for the IceCube Collaboration)}
{Proceedings of Science (PoS), ICRC2025; \href{https://arxiv.org/abs/2507.18525}{arXiv:2507.18525}}
\end{secbody}

\newpage
% === Conference Presentations ===
\section{Conference Presentations}
\conferenceitem{Apr 2026}{[Talk] Generative AI for IceCube Observatory}{103rd Spring Meeting of the Korean Physical Society, Daejeon, Republic of Korea}

\conferenceitem{Apr 2026}{[Poster] Mixture Density Networks for Fast Detector Response Simulation in SWGO}{103rd Spring Meeting of the Korean Physical Society, Daejeon, Republic of Korea}

\conferenceitem{Feb 2026}{[Poster] Investigation of IGZO-Based 3D Vertical Channel Transistors and Physics-Informed Neural Network Modeling for Next-Generation DRAM}{The 33rd Korean Conference on Semiconductors, Jeongseon, Republic of Korea}

\conferenceitem{Oct 2025}{[Poster] Diffusion and Flow Matching Model for IceCube Neutrino Event Simulation and Reconstruction}{102nd Fall Meeting of the Korean Physical Society, Gwangju, Republic of Korea}

\conferenceitem{Oct 2025}{[Poster] Generative AI for Cosmological Simulation and Manifold Representation}{KPS Undergraduate Physics Research Competition, 102nd Fall Meeting of the Korean Physical Society, Gwangju, Republic of Korea}

\conferenceitem{Jul 2025}{[Poster] Performance Study of the IceCube Upgrade Camera System}{39th International Cosmic Ray Conference (ICRC 2025), Geneva, Switzerland}

\conferenceitem{Jun 2025}{[Poster] Diffusion and Flow Matching Model for IceCube Neutrino Event Simulation and Reconstruction}{4th K-Neutrino Symposium, Seoul, Republic of Korea}

\conferenceitem{Apr 2025}{[Poster] Modeling the Relationship between PMT Response and Air-Shower Particles for SWGO}{101st Spring Meeting of the Korean Physical Society, Daejeon, Republic of Korea}

\conferenceitem{Oct 2024}{[Poster] Development of a Deep Learning Model for Analyzing Antarctic Ice Optical Properties}{KPS Undergraduate Physics Research Competition, 101st Fall Meeting of the Korean Physical Society, Yeosu, Republic of Korea}

\conferenceitem{Sep 2024}{[Talk] Deep Learning for Ice-Property Measurement Using Simulation Images}{IceCube Collaboration Meeting, Madison, Wisconsin, USA}

\conferenceitem{Jul 2024}{[Talk] Ice-Property Measurement Using a Camera System in the IceCube Upgrade}{3rd K-Neutrino Symposium, Gwangju, Republic of Korea}

\conferenceitem{Apr 2024}{[Poster] Study of Optical Properties in Antarctic Ice Using Simulated Images}{100th Spring Meeting of the Korean Physical Society, Daejeon, Republic of Korea}

\newpage
% === Honors & Awards ===
\section{Honors \& Awards}
\begin{secbody}
\awarditem{Best Presentation Award}{The 33rd Korean Conference on Semiconductors, 2026}{Awarded for poster presentation: ``Investigation of IGZO-Based 3D Vertical Channel Transistors and
Physics-Informed Neural Network Modeling for Next-Generation DRAM''}

\awarditem{Best Presentation Award}{4th K-Neutrino Symposium, 2025}{Awarded for poster presentation: ``Diffusion and Flow Matching Model for IceCube Neutrino Event Simulation and Reconstruction''}

\awarditem{Merit Award}{Korean Physical Society Undergraduate Physics Research Competition, 2024}{Recognized for research on CNN and Vision Transformer models for the IceCube Upgrade Camera System}
\end{secbody}

% === Skills ===
\section{Skills}
\begin{secbody}

\vspace{-1.5pt}
\begin{tabularx}{\linewidth}{@{}>{\bfseries}p{0.18\linewidth}X@{}}
Programming    & Python, NumPy, SciPy, Matplotlib \\[6pt]
Tools          & Linux, HTCondor, Git, Bash/Shell, HDF5 \\[6pt]
Frameworks     & PyTorch, MATLAB, JAX \\[6pt]
Architectures  & Transformer, Diffusion Transformer (DiT), Linformer, Vision Transformer (ViT), Swin Transformer, CNN, ResNet, U-Net, Perceiver \\[6pt]
Models         & VAE, Diffusion Models, Flow Models, Mixture Density Networks, Physics-Informed Neural Networks, Masked Autoencoder (MAE) \\[6pt]
Methods        & Optimal Transport Flow Matching, Rotary Positional Embedding (RoPE), Transformer Scaling Laws \\[6pt]
\end{tabularx}
\end{secbody}

% === Collaboration Memberships ===
\section{Past Collaboration Memberships}
\begin{secbody}

\noindent
\textbf{IceCube Collaboration} \hfill \textit{Aug 2023 -- Feb 2026} \\
{\small A neutrino detector embedded in Antarctic ice at the South Pole}

\vspace{8pt}

\noindent
\textbf{SWGO Collaboration} \hfill \textit{Jun 2024 -- Feb 2026} \\
{\small A next-generation ground-based gamma-ray observatory planned for the high-altitude Andes}
\end{secbody}

% === Languages ===
\section{Languages}
\begin{secbody}

\noindent
\textbf{Korean} -- Native \\
\textbf{English} -- Professional working proficiency (OPIc AL)
\end{secbody}

\end{document}
